\documentclass[12pt]{article}
\usepackage[utf8]{inputenc}
\usepackage[T1]{fontenc}
\usepackage{hyperref}
\title{Pana Language Resources for K-12 Curriculum Development}
\author{Compiled by Claude Sonnet 5.5 (claude-sonnet-5-5) for Andr\'as Kornai}
\date{}
\begin{document}
\maketitle

\section{Introduction}
Pana (ISO 639-3 pnz; Glottocode pana1293) is spoken in the Central African Republic, Chad and Cameroon. Wikipedia describes it as an Mbum language; Joshua Project calls it a Gur language, so sources disagree on classification. Speaker figures also vary (about 86,000 in one source, 2,980 in Cameroon per UNESCO). Listed dialects include Gonge, Man, Pana, Pani and Pondo. Very little was found.

\section{Primary Educational Resources}
No primers, textbooks or school programmes were found.

\section{Linguistic and Reference Materials}
\begin{itemize}
\item Pana language (Wikipedia): \url{https://en.wikipedia.org/wiki/Pana_language}
\item Pana/Pani numeral data (MPI-EVA): \url{https://lingweb.eva.mpg.de/channumerals/Pana-Pani.htm}
\item UNESCO World Atlas of Languages, Pana in Cameroon: \url{https://ar.wal.unesco.org/countries/cameroon/languages/pana}
\end{itemize}

\section{Religious and Community Resources}
\begin{itemize}
\item Joshua Project, Pana in Cameroon: \url{https://www.joshuaproject.net/index.php/people_groups/14285/CM}
\item Joshua Project, Pana in Central African Republic: \url{https://joshuaproject.net/index.php/people_groups/14285/CT}
\item Joshua Project, Pana in Chad: \url{https://joshuaproject.net/index.php/people_groups/14285/CD}
\item Joshua Project language page for pnz: \url{https://www.joshuaproject.net/index.php/languages/pnz}
\end{itemize}

\section{Digital and Audio Resources}
No audio, video or software resources were found.

\section{Curriculum Development Context}
There is no known orthography, literacy programme or school use found. The language spans three countries with different official school languages (French, Sango, Cameroonian French/English), so any work would start from documentation and community consultation.
\end{document}